\subsection{Comparison across families}
\label{sec:comparison}

\begin{table*}[!tb]
\centering\footnotesize\setlength\tabcolsep{3pt}\renewcommand{\arraystretch}{1.05}
\caption{What each method family licenses, the controls that license more, and where this corpus stops.
Column~2: what the artefact supports on its own, on held-out data above the study-stated reference
(notation of \S\ref{sec:prelim}, Eqs.~(\ref{eq:probe})--(\ref{eq:sae})). Column~3: tests that license a
strictly stronger reading, none a prerequisite (Table~\ref{tab:vcov}). Column~4: the gap that most limits
the family's evidence here, with the record counts of \S\ref{sec:findings}. Geometry is task-linked
geometry only (\S\ref{sec:models}).}
\label{tab:families}
\begin{tabular}{@{}L{2.6cm}L{4.55cm}L{4.15cm}L{4.35cm}@{}}
\toprule
Family & Claim licensed on its own & Stronger reading needs & Where this corpus stops \\
\midrule
\textbf{Structural and geometric} (task-linked) & shape of the space: gap, alignment, effective rank, CKA ($\Phi$ over $\{z_\lambda\}$); no concept & comparison model; random-initialised or shuffled baseline & comparison model \Vcomp{}/\Vcompn{}; reference distribution \Vncal{}/\Vncaln{} \\
\addlinespace[3pt]
\textbf{Supervised concept mapping and decodability} & $c$ recoverable from $z_\lambda$ by decoders in $\mathcal G$ (Eq.~\ref{eq:probe}) & control task or selectivity; shuffled labels; random encoder; fixed $\mathcal G$ across models & permutation null \Vperm{}/\Vpermn{}; random encoder \Vlrn{}/\Vlrnn{}; selectivity not yet measurable (\S\ref{sec:future}) \\
\quad\emph{zero-shot scoring} & $c$ recoverable by a fixed scorer ($\cos(z_v,z_t)$), once validated on held-out labels & prompt-set sensitivity; held-out validation & expert-validated concepts a minority; prompt sensitivity unreported \\
\quad\emph{concept activation vectors} & logit gradient-associated with a fitted direction (TCAV $S_{c,k}$), not recoverability & random-direction null; refit on shuffled labels & refit on shuffled labels: no record \\
\quad\emph{labelled similarity} & representation and concept distances correspond (RSA); no held-out value & permutation null over the labelled sets & permutation null: one record of five \\
\addlinespace[3pt]
\textbf{Feature decomposition and discovery} & a unit's activation tracks the named content ($D$, $a(x)$, Eq.~\ref{eq:sae}) & reconstruction $R^2$; validated names; cross-seed matching & semantic validation \Vsem{}/\Vsemn{}; cross-seed identity: none \\
\addlinespace[3pt]
\textbf{Attribution, localisation and tracing} & $c$ decodable at, or attended from, $(\ell,p,h)$; use unconfirmed & occlusion or patching confirmation; random-component comparison & input-space corroboration \Vinp{}/\Vinpn{}; component perturbation \Vint{}/\Vintn{} \\
\addlinespace[3pt]
\textbf{Representation intervention and control} (within the forward pass) & output sensitive to $z_\lambda$ ($\Delta B$, Eq.~\ref{eq:effect}); with matched controls, selectively to $v_c$ & matched control directions (\S\ref{sec:future}); equal-size random units & matched control \Vspca{}--\Vspc{}/\Vspcn{} targeted records; no dose--response test \\
\quad\emph{post-hoc transformation of stored embeddings} & changes what a refitted readout can do, not the model's computation & sham or alternative transformation; unrelated-attribute control & downstream evaluation only; no sham transformation \\
\quad\emph{parameter editing and unlearning} & behaviour changes as specified; no mediator unless measured & locality and retain-set checks; a measured mediator & no representation-level mediator identified \\
\midrule
\textbf{Intrinsic and concept-aligned designs} (tag) & label depends on $z_\lambda$ only through $\hat c$; prototypes and anchors mean what they are named & leakage test; random concept projections; intervention on $\hat c$; expert labels & concept-layer control \Vintrany{}/\Vintrn{}; leakage rarely measured \\
\bottomrule
\end{tabular}
\end{table*}
\begin{table*}[!tb]
\centering\footnotesize\setlength\tabcolsep{2.5pt}\renewcommand{\arraystretch}{1.05}
\caption{What each family costs and how far this corpus has taken it. \emph{Concept annotation}: whether
the design needs a concept vocabulary, and how much labelling. \emph{Loci}: where the family has been
applied to medical models (parentheses: reached only in general-domain work). \emph{Compute}: beyond a
forward pass. \emph{Readable}: whether a clinician can read the artefact directly. \emph{Stud./rec.}:
core studies whose primary family is this one / records of this family (for decoding,
\Ndecstudnb{}/\Ndecrecnb{} without the \Nbenchdec{} frozen-feature benchmarks). \emph{Groups}: distinct
first authors among those records. \emph{Base.}: records comparing a medical with a general model on the
same data. \emph{Multi.}: more than one model or dataset. The structural row is task-linked geometry only
(\S\ref{sec:models}). Counts come from the coding sheet, the other cells are the authors' reading, and
the \Nintrinsic{} intrinsic-design records are counted in their analytic family (thirteen decoding, one
discovery) and listed again in the tag row.}
\label{tab:famreq}
\begin{tabular}{@{}L{1.86cm}L{1.82cm}L{1.76cm}L{1.33cm}L{1.53cm}C{0.49cm}C{0.89cm}C{0.79cm}C{0.61cm}C{0.61cm}L{2.55cm}@{}}
\toprule
Family & Concept annotation & Loci (medical) & Compute & Readable & 1st yr & Stud./ rec. & Groups & Base. & Multi. & Main failure mode \\
\midrule
Structural and geometric (task-linked) & none & Enc, Txt, Joint, Dec, (Conn) & low & no & 2022 & 21/30 & 27 & 11 & 29 & geometric quantity read as clinical meaning \\
Supervised concept mapping and decodability & required; high (labels) & Enc, Joint & low & partly (scores yes, probes no) & 2022 & 75/80 & 76 & 31 & 72 & decodability read as use \\
Feature decomposition and discovery & none or weak (post-hoc naming) & Enc, Dec, (Conn) & high (activations at scale, training) & yes, once named and validated & 2024 & 20/25 & 25 & 2 & 14 & unstable features and unvalidated names \\
Attribution, localisation and tracing & optional; low--moderate & Enc, Joint, Dec, Path, (Pre) & moderate & partly (maps, often unfaithful) & 2022 & 9/11 & 11 & 2 & 8 & localisation read as causal confirmation \\
Representation intervention and control & required for concept-\hspace{0pt}specific claims; moderate (exemplars) & Enc, Joint, Dec, (Txt, Conn) & mod.--high (repeated generation) & indirectly (behaviour change) & 2023 & 10/21 & 20 & 5 & 18 & uncontrolled, input-\hspace{0pt}dependent effects \\
\midrule
Intrinsic and concept-aligned designs (tag) & required; high & Enc (concept layer), Joint & low--moderate & yes & 2023 & 14 rec. & 13 & 7 & 13 & leakage; concept invalidity \\
\bottomrule
\end{tabular}
\end{table*}

Table~\ref{tab:families} sets out what each family reports and what its artefact licenses;
Table~\ref{tab:famreq} gives what each costs and how far this corpus has taken it, placing the medical
evidence at the observational end (Fig.~\ref{fig:corpus}b). Within this corpus, and under its family-specific eligibility rules, decoding is the only family with a broad base of distinct first authors (76); discovery is two years old and intervention three, and they have 25 and 20 distinct first authors, a proxy for independent groups rather than a measure of the families' maturity in the field. Two asymmetries in Table~\ref{tab:famreq} are easily
missed: intervention gives the claims of greatest practical interest at the highest cost and with the fewest
controlled results, while the cheapest family addresses availability, often mistaken for use; and readability
and validity pull apart, since discovery and intrinsic designs return the artefacts a clinician can read yet
rest on an unvalidated naming or leakage step.

\textbf{Combining families.} Three pipelines recur. \emph{Discover, then intervene} uses residual-stream
dictionary directions for steering or
patching~\cite{bouzid2025insights,nooralahzadeh2026universal,sadanandan2026mechanistically}, carrying the
caveats of \S\ref{sec:discovery} into the intervention. \emph{Trace, then intervene} exploits an
observational localisation through a training-time change, so far only in a general CLIP text tower on
radiology captions~\cite{abbasi2026layerspecific}. \emph{Decode, then erase, residualise or ablate} probes
for availability, then tests whether a downstream prediction depends on the
concept~\cite{kim2026encoding,nooralahzadeh2026sparse,weber2023posthoc}, answering the shortcut question of
\S\ref{sec:applications} most directly. Of its three core instances the two that test a downstream
dependence point in opposite directions, since
residualising a demographic attribute leaves a refitted finding probe unchanged~\cite{kim2026encoding}$^{\dagger}$
whereas knocking out the roughly ten channels that decode a finding from a CT encoder collapses its
score~\cite{nooralahzadeh2026sparse}; the third orthogonalises chest-radiograph embeddings and reports a
downstream change within noise~\cite{weber2023posthoc}.

Reading one family's result as an answer to another family's question is the commonest inferential error in
this literature, and the evidence-type tag makes it visible.

